\coursepart{1}{(MATH101) Calculus I}
\ifCASolutions\else
{\noindent\small\itshape Problems within each course are arranged chronologically by year and semester, then by institution in the order \POSTECH{}, \KAIST{}, and \SNU{}, while preserving the original problem order within each source.\par}
\vspace{0.22cm}
\fi

% 2019 Spring POSTECH Calculus I First Exam, Problem 1
\begin{problem}[2019 Spring \POSTECH{} Calculus I]
Evaluate
\[
\int \sin^3 x\,\dd x.
\]
\end{problem}

% 2019 Spring POSTECH Calculus I First Exam, Problem 2
\begin{problem}[2019 Spring \POSTECH{} Calculus I]
Evaluate
\[
\int \frac{x}{(x+2)(x+1)^2}\,\dd x.
\]
\end{problem}

% 2019 Spring POSTECH Calculus I First Exam, Problem 3
\begin{problem}[2019 Spring \POSTECH{} Calculus I]
Evaluate
\[
\int \frac{1}{5-3\cos x}\,\dd x.
\]
\end{problem}

% 2019 Spring POSTECH Calculus I First Exam, Problem 4
\begin{problem}[2019 Spring \POSTECH{} Calculus I]
Test the convergence of the series
\[
\sum_{n=1}^{\infty}
\frac{5^{3n^2}-6^n}{5^{3n^2}+7^n\sqrt n}.
\]
\end{problem}

% 2019 Spring POSTECH Calculus I First Exam, Problem 5
\begin{problem}[2019 Spring \POSTECH{} Calculus I]
Test the convergence of the series
\[
\sum_{n=2}^{\infty}
\frac{\cos^n n}{2^n-\sin(n^3)}.
\]
\end{problem}

% 2019 Spring POSTECH Calculus I Midterm, Problem 1
\begin{problem}[2019 Spring \POSTECH{} Calculus I]
Determine the values of the real number $\alpha$ for which the series
\[
\sum_{n=1}^{\infty}\frac{(\ln n)^\alpha}{n}
\]
converges.
\end{problem}

% 2019 Spring POSTECH Calculus I Midterm, Problem 2
\begin{problem}[2019 Spring \POSTECH{} Calculus I]
Evaluate
\[
\int \frac{1}{(1+x^2)^2}\,\dd x.
\]
\end{problem}

% 2019 Spring POSTECH Calculus I Midterm, Problem 3
\begin{problem}[2019 Spring \POSTECH{} Calculus I]
Compute the radius of convergence of
\[
f(x)=\sum_{n=0}^{\infty}(-1)^n x^{2n}
\]
and the interval of convergence of
\[
g(x)=\int_0^x f(y)\,\dd y,
\]
respectively.
\end{problem}

% 2019 Spring POSTECH Calculus I Midterm, Problem 4
\begin{problem}[2019 Spring \POSTECH{} Calculus I]
Find the Taylor series of
\[
f(x)=\ln\!\left(x+\sqrt{1+x^2}\right)
\]
about $x=0$, and also find its radius of convergence.
\end{problem}

% 2019 Spring POSTECH Calculus I Midterm, Problem 5
\begin{problem}[2019 Spring \POSTECH{} Calculus I]
Find the radius of convergence of
\[
\sum_{n=1}^{\infty}a_nx^n,
\]
where
\[
\frac{1}{\sqrt n}\le a_n\le\sqrt n
\]
for every integer $n$.
\end{problem}

% 2019 Spring POSTECH Calculus I Midterm, Problem 6
\begin{problem}[2019 Spring \POSTECH{} Calculus I]
Find the function of the form
\[
f(x)=\sum_{n=0}^{\infty}a_n\frac{x^n}{n!}
\]
which satisfies
\[
\frac{\dd^3}{\dd x^3}f(x)=-f(x),
\qquad
f(0)=1,
\qquad
f'(0)=0,
\qquad
f''(0)=2.
\]
\end{problem}

% 2019 Spring POSTECH Calculus I Midterm, Problem 7
\begin{problem}[2019 Spring \POSTECH{} Calculus I]
\begin{problemsubparts}
\item Prove that
\[
\cos x=\sum_{n=0}^{\infty}(-1)^n\frac{x^{2n}}{(2n)!}.
\]
\item Find the limit
\[
\lim_{x\to0}
\frac{\cos(x^6)-1+\frac{1}{2!}x^{12}-\frac{x^{24}}{4!}}{x^{36}}.
\]
\end{problemsubparts}
\end{problem}

% 2019 Spring POSTECH Calculus I Midterm, Problem 8
\begin{problem}[2019 Spring \POSTECH{} Calculus I]
Derive the Taylor series of
\[
f(x)=\ln(2+3x)
\]
about $x=0$, and determine the set of points $x$ at which its Taylor series $T_0f(x)$ converges to $f(x)$.
\end{problem}

% 2019 Spring POSTECH Calculus I Midterm, Problem 9
\begin{problem}[2019 Spring \POSTECH{} Calculus I]
Find an approximate value of
\[
\int_0^1 e^{-x^2}\,\dd x
\]
with allowed error not more than $0.01$. Justify your answer.
\end{problem}

% 2019 Spring POSTECH Calculus I Midterm, Problem 10
\begin{problem}[2019 Spring \POSTECH{} Calculus I]
Assume that $a_n$, for every $n=0,1,2,\ldots$, is an integer between $1$ and $10^8$. Find the radius of convergence of the series
\[
\sum_{n=0}^{\infty}\bigl(a_n+1000\cos^2 n\bigr)x^n.
\]
\end{problem}

% 2019 Spring POSTECH Calculus I Final, Problem 1
\begin{problem}[2019 Spring \POSTECH{} Calculus I]
Determine whether
\[
\sum_{n=5}^{\infty}
\frac{(n+1)^n}{n^{n+1}(\ln n)^2}
\]
converges.
\end{problem}

% 2019 Spring POSTECH Calculus I Final, Problem 2
\begin{problem}[2019 Spring \POSTECH{} Calculus I]
Determine whether the sequence
\[
A_n
=2+\frac{2}{3}+\frac{2}{5}+\cdots+\frac{2}{2n-1}-\ln(2n+1)
\]
converges as $n$ tends to infinity.
\end{problem}

% 2019 Spring POSTECH Calculus I Final, Problem 3
\begin{problem}[2019 Spring \POSTECH{} Calculus I]
Evaluate
\[
\int x^2\arcsin x\,\dd x.
\]
\end{problem}

% 2019 Spring POSTECH Calculus I Final, Problem 4
\begin{problem}[2019 Spring \POSTECH{} Calculus I]
Evaluate
\[
\int \frac{x}{(x+1)^2(x+2)}\,\dd x.
\]
\end{problem}

% 2019 Spring POSTECH Calculus I Final, Problem 5
\begin{problem}[2019 Spring \POSTECH{} Calculus I]
Find the radius of convergence of the series
\[
\sum_{n=1}^{\infty}
\frac{\sin\!\left(\frac{n\pi}{2}\right)}{2^n}x^n.
\]
\end{problem}

% 2019 Spring POSTECH Calculus I Final, Problem 6
\begin{problem}[2019 Spring \POSTECH{} Calculus I]
Find an approximate value of
\[
\int_0^1\sin(x^2)\,\dd x
\]
with allowed error not more than $0.001$. Justify your answer.
\end{problem}

% 2019 Spring POSTECH Calculus I Final, Problem 7
\begin{problem}[2019 Spring \POSTECH{} Calculus I]
Solve for the function $y$ from
\[
y'+y=e^{-x+1},
\qquad
y(0)=e.
\]
\end{problem}

% 2019 Spring POSTECH Calculus I Final, Problem 8
\begin{problem}[2019 Spring \POSTECH{} Calculus I]
Let $f\colon\R\to\R$ be defined by
\[
f(x)=
\begin{cases}
x^2\sin\!\left(\dfrac{1}{x}\right), & x\ne0,\\[0.6em]
0, & x=0.
\end{cases}
\]
Show that $f$ is differentiable at $x=0$.
\end{problem}

% 2019 Spring POSTECH Calculus I Final, Problem 9
\begin{problem}[2019 Spring \POSTECH{} Calculus I]
Assume that the solution to the differential equation
\[
y'=y^2+1,
\qquad
y(0)=0,
\]
is a power series of the form
\[
\sum_{n=0}^{\infty}a_nx^n
\]
that converges for every $x$ with $|x|<r$ for some positive number $r$. Find the first two nonzero terms of the Taylor series of the solution function centered at $0$.
\end{problem}

% 2019 Spring POSTECH Calculus I Final, Problem 10
\begin{problem}[2019 Spring \POSTECH{} Calculus I]
The functions $f$, $g$, and $h$ are defined on the interval $[-1,1]$ and satisfy
\[
f(x)\le g(x)\le h(x)
\]
for every $x\in[-1,1]$. Assume also that
\[
\lim_{x\to0}f(x)=3=\lim_{x\to0}h(x).
\]
Using the $\varepsilon$-$\delta$ method, that is, the definition of limit, show that
\[
\lim_{x\to0}g(x)=3.
\]
\end{problem}

% 2021 Spring POSTECH Calculus I Midterm, Problem 1
\begin{problem}[2021 Spring \POSTECH{} Calculus I]
Let
\[
f(x)=(1+x^2)^{72}.
\]
Evaluate $f^{(70)}(0)$.
\end{problem}

% 2021 Spring POSTECH Calculus I Midterm, Problem 2
\begin{problem}[2021 Spring \POSTECH{} Calculus I]
Evaluate
\[
\int \frac{x}{(x+2)(x+1)^2}\,\dd x.
\]
\end{problem}

% 2021 Spring POSTECH Calculus I Midterm, Problem 3
\begin{problem}[2021 Spring \POSTECH{} Calculus I]
Suppose
\[
e^{-x^2}\sin 2x=\sum_{n=0}^{\infty}a_nx^n
\]
for all $x\in\R$. Find $a_{100}$.
\end{problem}

% 2021 Spring POSTECH Calculus I Midterm, Problem 4
\begin{problem}[2021 Spring \POSTECH{} Calculus I]
Find the interval of convergence of the power series
\[
\sum_{n=1}^{\infty}(-1)^n\frac{x^{n+1}}{n+1}.
\]
\end{problem}

% 2021 Spring POSTECH Calculus I Midterm, Problem 5
\begin{problem}[2021 Spring \POSTECH{} Calculus I]
Find the interval of convergence of the power series
\[
\sum_{n=1}^{\infty}
\frac{2^n+(n-1)!}{n!}x^n.
\]
\end{problem}

% 2021 Spring POSTECH Calculus I Midterm, Problem 6
\begin{problem}[2021 Spring \POSTECH{} Calculus I]
Show that
\[
\lim_{n\to\infty}\frac{100^n}{n!}=0.
\]
\end{problem}

% 2021 Spring POSTECH Calculus I Midterm, Problem 7
\begin{problem}[2021 Spring \POSTECH{} Calculus I]
Consider an approximation of
\[
\int_0^{0.1}e^{-x^2}\,\dd x
\]
using the Taylor polynomial of $e^{-x^2}$ with two nonzero terms. Show that the error of this approximation does not exceed $10^{-6}$.
\end{problem}

% 2021 Spring POSTECH Calculus I Midterm, Problem 8
\begin{problem}[2021 Spring \POSTECH{} Calculus I]
\begin{problemsubparts}
\item Find the Taylor series, denoted by $T_0f(x)$, for $f(x)=\arctan x$ at $a=0$.
\item Prove that
\[
T_0f(x)=\arctan x
\]
for $-1\le x\le1$. You must also show that the equality does not hold for $|x|>1$.

\textit{Hint.} You may want to use the fact that
\[
\frac{|f^{(n)}(x)|}{n!}\le\frac1n
\]
for $|x|\le1$.
\end{problemsubparts}
\end{problem}

% 2021 Spring KAIST Calculus I Midterm, Problem 1
\begin{problem}[2021 Spring \KAIST{} Calculus I]
Does each of the following infinite series converge? Mark Yes or No. You do not need to justify your answers.
\begin{multicols}{2}
\begin{problemchoices}
\item $\displaystyle\sum_{n=1}^{\infty}\frac{n!}{n^n}$
\item $\displaystyle\sum_{n=1}^{\infty}\frac{e^{\sqrt n}}{n\sqrt n}$
\item $\displaystyle\sum_{n=1}^{\infty}\left(1-\frac1n\right)^n$
\item $\displaystyle\sum_{n=1}^{\infty}(-1)^n\left(1-\frac1n\right)^{n\ln n}$
\item $\displaystyle\sum_{n=1}^{\infty}\left(1-\frac1n\right)^{2n\ln n}$
\item $\displaystyle\sum_{n=1}^{\infty}\frac{1}{(n+1)\ln(n+1)}$
\item $\displaystyle\sum_{n=1}^{\infty}\frac{1}{(n+1)\ln(n+1)\ln\ln(n+1)}$
\item $\displaystyle\sum_{n=1}^{\infty}\frac{1}{(n+1)(\ln(n+1))^2}$
\item $\displaystyle\sum_{n=1}^{\infty}\frac{1}{\tan\!\left(\arccos\!\left(\frac{1}{n+1}\right)\right)}$
\item $\displaystyle\sum_{n=1}^{\infty}\frac{\sin(1/n)}{n^\alpha},\qquad \alpha>0$
\end{problemchoices}
\end{multicols}
\end{problem}

% 2021 Spring KAIST Calculus I Midterm, Problem 2
\begin{problem}[2021 Spring \KAIST{} Calculus I]
Find the area of the portion of the ellipse
\[
9x^2+16y^2=144
\]
between the $y$-axis and $x=2$.
\end{problem}

% 2021 Spring KAIST Calculus I Midterm, Problem 3
\begin{problem}[2021 Spring \KAIST{} Calculus I]
Find the antiderivative of
\[
\frac{1}{(x^2+x+1)(x^2+x+2)(x^2+x+3)}.
\]
\end{problem}

% 2021 Spring KAIST Calculus I Midterm, Problem 4
\begin{problem}[2021 Spring \KAIST{} Calculus I]
Does the following improper integral converge?
\[
\int_0^\infty \frac{\sin^3 x}{x}\,\dd x.
\]
\end{problem}

% 2021 Spring KAIST Calculus I Midterm, Problem 5
\begin{problem}[2021 Spring \KAIST{} Calculus I]
Prove or disprove:
\begin{problemchoices}
\item If $\displaystyle\sum_{n=1}^{\infty}a_n$ converges, then $\displaystyle\sum_{n=1}^{\infty}a_n^2$ also converges.
\item If both $\displaystyle\sum_{n=1}^{\infty}a_n$ and $\displaystyle\sum_{n=1}^{\infty}|b_n|$ converge, then $\displaystyle\sum_{n=1}^{\infty}a_nb_n$ also converges.
\item Assume $a_n\ne0$ for all $n$. If $\displaystyle\lim_{n\to\infty}a_n=0$ and $\displaystyle\sum_{n=1}^{\infty}|\sin(a_n)|$ converges, then $\displaystyle\sum_{n=1}^{\infty}|a_n|$ converges.
\end{problemchoices}
\end{problem}

% 2021 Spring KAIST Calculus I Midterm, Problem 6
\begin{problem}[2021 Spring \KAIST{} Calculus I]
Evaluate
\[
\int_1^\infty\left(2\sqrt{x^2-1}-2x+\frac1x\right)\dd x.
\]
\end{problem}

% 2021 Spring KAIST Calculus I Midterm, Problem 7
\begin{problem}[2021 Spring \KAIST{} Calculus I]
Let $f$ be a continuously differentiable (i.e., $f(x)$ is differentiable and $f'(x)$ is continuous) one-to-one function defined on $[a,b]$. Assume that $f'(x)\ne0$ on $[a,b]$. Compute the following sum of two integrals:
\[
\int_{f(a)}^{f(b)}f^{-1}(y)\,\dd y+\int_a^b f(x)\,\dd x.
\]
\end{problem}

% 2021 Spring KAIST Calculus I Midterm, Problem 8
\begin{problem}[2021 Spring \KAIST{} Calculus I]
Let $n$ be a positive integer and denote the following improper integral by $a_n$:
\[
a_n=\int_0^1(\ln x)^n\,\dd x.
\]
\begin{problemsubparts}
\item Show that the improper integral $a_n$ is convergent. Compute the value of $a_n$.
\item Is the series
\[
\sum_{n=1}^{\infty}\frac1{a_n}
\]
absolutely convergent, conditionally convergent, or divergent? Explain the reason using appropriate tests.
\item Is the series
\[
\sum_{n=2}^{\infty}(-1)^n\frac1{\ln|a_n|}
\]
absolutely convergent, conditionally convergent, or divergent? Explain the reason using appropriate tests.
\end{problemsubparts}
\end{problem}

% 2021 Spring KAIST Calculus I Final, Problem 3
\begin{problem}[2021 Spring \KAIST{} Calculus I]
Show that the power series
\[
\sum_{n=2}^{\infty}n(n-1)x^n,
\qquad |x|<1,
\]
can be expressed as
\[
\frac{x^2}{(1-x)^3}.
\]
\end{problem}

% 2021 Spring KAIST Calculus I Final, Problem 4
\Needspace{18\baselineskip}
\begin{problem}[2021 Spring \KAIST{} Calculus I]
Consider the standard ellipse $C$,
\[
\frac{x^2}{a^2}+\frac{y^2}{b^2}=1,
\qquad a>b>0.
\]
Denote by $e$ the eccentricity of $C$.
\begin{problemsubparts}
\item Find an upper bound of the function
\[
f(t)=\sqrt{1-e^2\sin^2 t}
\]
using the first two terms of its binomial series expansion.
\item Show that the perimeter (total arc length) of $C$ is less than
\[
2\pi a\left(1-\frac{e^2}{4}\right).
\]
\end{problemsubparts}
\end{problem}

% 2021 Spring KAIST Calculus I Final, Problem 6
\begin{problem}[2021 Spring \KAIST{} Calculus I]
Consider the power series
\[
\sum_{n=1}^{\infty}(-1)^n\frac{(x-1)^{2n+1}}{n2^n}.
\]
\begin{problemsubparts}
\item Find its interval of convergence.
\item At each point in the interval of convergence, discuss its mode of convergence as absolute or conditional convergence together with proper reasoning.
\end{problemsubparts}
\end{problem}

% 2021 Fall KAIST Calculus I Midterm, Problem 1
\begin{problem}[2021 Fall \KAIST{} Calculus I]
Is each of the following statements true or false? Mark T or F. You do not need to justify your answers.
\begin{problemchoices}
\item $f=O(g)$ implies $f=o(g)$ for functions $f$ and $g$ that are positive for $x$ sufficiently large.
\item $x(\ln x)^{2022}=o(x^{1.001})$.
\item If $f'(a)\ne0$, then there exists a $\delta>0$ such that $f(x)\ne f(a)$ for all $x\ne a$ with $a-\delta<x<a+\delta$.
\item $\displaystyle\int_2^\infty\frac{1-2^{-x}}{1+\sqrt x}\,\dd x$ converges.
\item $\displaystyle\int_1^6\frac{1}{x-4}\,\dd x$ converges.
\item $\displaystyle\int_2^\infty\frac{1}{\ln x}\,\dd x$ converges.
\item $\displaystyle\int_1^2\frac{1}{x\ln x}\,\dd x$ converges.
\item $\displaystyle\int_{-\infty}^{\infty}\frac{1}{e^x+e^{-x}}\,\dd x$ converges.
\item $\displaystyle\int_{-\infty}^{\infty}\frac{2x}{x^2+1}\,\dd x$ converges.
\item $\displaystyle\lim_{x\to\infty}\frac{\ln(x+10^{2022})}{\ln x}=1$.
\end{problemchoices}
\end{problem}

% 2021 Fall KAIST Calculus I Midterm, Problem 2
\begin{problem}[2021 Fall \KAIST{} Calculus I]
Evaluate
\[
\int_1^\infty\left(2\sqrt{x^2-1}-2x+\frac1x\right)\dd x.
\]
\end{problem}

% 2021 Fall KAIST Calculus I Midterm, Problem 3
\begin{problem}[2021 Fall \KAIST{} Calculus I]
Let $f$ be a continuously differentiable (i.e., $f(x)$ is differentiable and $f'(x)$ is continuous) one-to-one function defined on $[a,b]$. Assume that $f'(x)\ne0$ on $[a,b]$. Compute the following sum of two integrals:
\[
\int_{f(a)}^{f(b)}f^{-1}(y)\,\dd y+\int_a^b f(x)\,\dd x.
\]
\end{problem}

% 2021 Fall KAIST Calculus I Midterm, Problem 4
\begin{problem}[2021 Fall \KAIST{} Calculus I]
\begin{problemsubparts}
\item Derive $\displaystyle\frac{\dd}{\dd x}(\arctanh x)$ for $|x|<1$.
\item Find $\displaystyle\int x\arctanh x\,\dd x$.
\end{problemsubparts}
\end{problem}

% 2021 Fall KAIST Calculus I Midterm, Problem 5
\begin{problem}[2021 Fall \KAIST{} Calculus I]
Let
\[
\Gamma(n)=\int_0^\infty x^{n-1}e^{-x}\,\dd x
\]
for a positive integer $n$.
\begin{problemsubparts}
\item Use a comparison test to show that $\Gamma(n)$ converges.
\item Show that $\Gamma(n)=(n-1)!$.
\end{problemsubparts}
\end{problem}

% 2021 Fall KAIST Calculus I Midterm, Problem 6
\begin{problem}[2021 Fall \KAIST{} Calculus I]
Let $w(t)$ be a differentiable function and let $\kappa$ and $\rho$ be positive constants.
\begin{problemsubparts}
\item Solve
\[
\frac{\dd}{\dd t}w(t)+\kappa w(t)-\rho=0
\qquad\text{or}\qquad
w'+\kappa w-\rho=0
\]
with the initial condition $w(0)=v_0$.
\item By setting $w(t)=\displaystyle\frac{\dd}{\dd t}y(t)$, solve
\[
\frac{\dd^2}{\dd t^2}y(t)+\kappa\frac{\dd}{\dd t}y(t)-\rho=0
\qquad\text{or}\qquad
y''+\kappa y'-\rho=0
\]
with $y(0)=0$ and
\[
\left.\frac{\dd}{\dd t}y(t)\right|_{t=0}=v_0.
\]
\end{problemsubparts}
\end{problem}

% 2021 Fall KAIST Calculus I Midterm, Problem 7
\begin{problem}[2021 Fall \KAIST{} Calculus I]
\begin{problemchoices}
\item Solve the equation
\[
xy'+2y=\frac{8x}{(x-1)(x+1)(x+2)}
\]
for $x>1$.
\item Solve the initial value problem
\[
y'=2x\sqrt{1+y^2},
\qquad y(0)=0.
\]
\end{problemchoices}
\end{problem}

% 2021 Fall KAIST Calculus I Final, Problem 3
\begin{problem}[2021 Fall \KAIST{} Calculus I]
Determine the radius and interval of convergence of the following series.
\begin{problemchoices}
\item $\displaystyle\sum_{n=1}^{\infty}\frac1n x^n$
\item $\displaystyle\sum_{n=1}^{\infty}\frac1{n!}x^n$
\item $\displaystyle\sum_{n=0}^{\infty}(-1)^n x^{2^n}$
\end{problemchoices}
\end{problem}

% 2021 Fall KAIST Calculus I Final, Problem 4
\begin{problem}[2021 Fall \KAIST{} Calculus I]
Consider the power series
\[
\sum_{n=1}^{\infty}(-1)^n\frac{(x-1)^{2n+1}}{n2^n}.
\]
\begin{problemsubparts}
\item Find its interval of convergence.
\item At each point in the interval of convergence, discuss its mode of convergence as absolute or conditional convergence together with proper reasoning.
\end{problemsubparts}
\end{problem}

% 2022 Spring POSTECH Calculus I Midterm, Problem 1
\begin{problem}[2022 Spring \POSTECH{} Calculus I]
Prove that the sequence $(a_n)$ defined by
\[
a_n=\frac{1}{n^2},
\qquad n\ge1,
\]
is a Cauchy sequence.
\end{problem}

% 2022 Spring POSTECH Calculus I Midterm, Problem 2
\begin{problem}[2022 Spring \POSTECH{} Calculus I]
A sequence $(s_n)$ is defined by
\[
s_1=7,
\qquad
s_n=\frac12\left(s_{n-1}+\frac{7}{s_{n-1}}\right),
\qquad n\ge2.
\]
\begin{problemsubparts}
\item Show that $(s_n)$ is bounded below.

\textit{Hint.} Use the arithmetic--geometric mean inequality.
\item Show that $(s_n)$ is decreasing.
\end{problemsubparts}
\end{problem}

% 2022 Spring POSTECH Calculus I Midterm, Problem 3
\begin{problem}[2022 Spring \POSTECH{} Calculus I]
Show that
\[
f_n(x)=\sum_{k=0}^{n-1}x^k
\]
converges uniformly on $[-c,c]$, where $0<c<1$.
\end{problem}

% 2022 Spring POSTECH Calculus I Midterm, Problem 4
\begin{problem}[2022 Spring \POSTECH{} Calculus I]
The function
\[
f(x)=\frac{x^3-9x}{x^2+3x}
\]
is defined on $[1,10]$.
\begin{problemsubparts}
\item For $\varepsilon=10^{-3}$, find $\delta>0$ such that if $|x-5|<\delta$, then
\[
|f(x)-f(5)|<\varepsilon.
\]
\item For any $\varepsilon>0$, find $\delta>0$ such that if $|x-5|<\delta$, then
\[
|f(x)-f(5)|<\varepsilon.
\]
\end{problemsubparts}
\end{problem}

% 2022 Spring POSTECH Calculus I Midterm, Problem 5
\begin{problem}[2022 Spring \POSTECH{} Calculus I]
Let $f\colon\R\to\R$ be defined by
\[
f(x)=
\begin{cases}
x^2\sin\!\left(\dfrac1x\right), & x\ne0,\\[0.5em]
0, & x=0.
\end{cases}
\]
\begin{problemsubparts}
\item Show that $f$ is differentiable at $0$, and find $f'(0)$.
\item Is $f'$ continuous at $x=0$? Justify your answer.
\end{problemsubparts}
\end{problem}

% 2022 Spring POSTECH Calculus I Midterm, Problem 6
\begin{problem}[2022 Spring \POSTECH{} Calculus I]
Suppose that at time $t=1\,\mathrm{h}$, the population of a bacterial colony is $100$, and the rate of population growth per hour is $1.5$ times the size of the population. Find the population at time $t=3\,\mathrm{h}$.
\end{problem}

% 2022 Spring POSTECH Calculus I Midterm, Problem 7
\begin{problem}[2022 Spring \POSTECH{} Calculus I]
Suppose a sequence $(a_n)$ converges to $a$. Use and justify the following steps to prove that the sequence is bounded.
\begin{problemsubparts}
\item There is a number $N$ such that, for all $n>N$,
\[
|a_n-a|<1.
\]
\item For all $n>N$,
\[
|a_n|\le |a|+1.
\]
\item Let $\alpha$ be the largest of the numbers
\[
|a_1|,|a_2|,\ldots,|a_N|,|a|+1.
\]
Then $|a_k|\le\alpha$ for $k=1,2,\ldots$.
\item The sequence $(a_n)$ is bounded.
\end{problemsubparts}
\end{problem}

% 2022 Spring POSTECH Calculus I Midterm, Problem 8
\begin{problem}[2022 Spring \POSTECH{} Calculus I]
Suppose that a series $\sum a_n$ is convergent and $a_n\ge0$. Find the interval of convergence of
\[
\sum_{n=0}^{\infty}
\left(\frac{1+\sin a_n}{2}\right)^n x^n.
\]
\end{problem}

% 2022 Spring POSTECH Calculus I Midterm, Problem 9
\begin{problem}[2022 Spring \POSTECH{} Calculus I]
\begin{problemsubparts}
\item Let
\[
f(x)=\sum_{i=1}^{n}(a_i x+b_i)^2,
\qquad x\in(-\infty,\infty),
\]
where $a_i$ and $b_i$ are real numbers. Using the quadratic function $f(x)$, derive the Cauchy--Schwarz inequality
\[
\sum_{j=1}^{n}a_jb_j
\le
\sqrt{\sum_{j=1}^{n}a_j^2}\,
\sqrt{\sum_{j=1}^{n}b_j^2}.
\]
\item Use the Cauchy--Schwarz inequality to show that
\[
\sum_{n=0}^{\infty}a_nb_n
\]
converges absolutely if both
\[
\sum_{n=0}^{\infty}a_n^2
\qquad\text{and}\qquad
\sum_{n=0}^{\infty}b_n^2
\]
converge.
\end{problemsubparts}
\end{problem}

% 2022 Spring POSTECH Calculus I Midterm, Problem 10
\begin{problem}[2022 Spring \POSTECH{} Calculus I]
Let
\[
f(x)=\sinh x=\frac{e^x-e^{-x}}{2},
\qquad x\in(-\infty,\infty).
\]
\begin{problemsubparts}
\item Show that $f$ is invertible.
\item Derive a formula for $\arcsinh y$.
\item Find the domain and range of $\arcsinh y$.
\item Find the derivative of $\arcsinh y$.
\end{problemsubparts}
\end{problem}

% 2022 Spring POSTECH Calculus I Midterm, Problem 11
\begin{problem}[2022 Spring \POSTECH{} Calculus I]
Use the definition of the derivative to show the following: if $f$ is differentiable at $x$ and $f(x)\ne0$, then $1/f$ is differentiable at $x$, and
\[
\left(\frac1f\right)'(x)
=-\frac{f'(x)}{(f(x))^2}.
\]
\end{problem}

% 2022 Spring POSTECH Calculus I Final, Problem 1
\begin{problem}[2022 Spring \POSTECH{} Calculus I]
Let $V(h)$ be the volume of the solid obtained by rotating the region bounded by
\[
y=\frac{1}{x},
\qquad y=1,
\qquad y=h,
\qquad x=0
\]
about the $y$-axis. Find
\[
\lim_{h\to\infty}V(h).
\]
\end{problem}

% 2022 Spring POSTECH Calculus I Final, Problem 2
\begin{problem}[2022 Spring \POSTECH{} Calculus I]
Let $f(x)=\sin x$, and let $T(x)$ be the Taylor series of $f$ centered at $x=\pi/6$.
\begin{problemsubparts}
\item Find the first four nonzero terms of $T(x)$.
\item Show that $f(x)=T(x)$ for all $x\in\R$.
\end{problemsubparts}
\end{problem}

% 2022 Spring POSTECH Calculus I Final, Problem 3
\begin{problem}[2022 Spring \POSTECH{} Calculus I]
Let
\[
f_n(x)=n^3x(1-x^2)^n,
\qquad x\in[0,1].
\]
\begin{problemsubparts}
\item Show that $f_n(x)$ converges to zero pointwise.
\item Calculate
\[
\int_0^1 f_n(x)\,\dd x.
\]
\item Show that
\[
\lim_{n\to\infty}\int_0^1 f_n(x)\,\dd x
\ne
\int_0^1\lim_{n\to\infty}f_n(x)\,\dd x.
\]
\end{problemsubparts}
\end{problem}

% 2022 Spring POSTECH Calculus I Final, Problem 4
\begin{problem}[2022 Spring \POSTECH{} Calculus I]
Determine whether the series
\[
\sum_{n=1}^{\infty} n^3e^{-n^4}
\]
converges. Justify your answer.
\end{problem}

% 2022 Spring POSTECH Calculus I Final, Problem 5
\begin{problem}[2022 Spring \POSTECH{} Calculus I]
Let $f\colon[a,b]\to\R$ be a function.
\begin{problemsubparts}
\item If $f$ is continuously differentiable, show that there exists a continuous function $g\colon[a,b]\to\R$ such that
\[
\int_a^x g(t)\,\dd t=f(x)-f(a)
\]
for all $x\in[a,b]$.
\item Conversely, if there exists a continuous function $g\colon[a,b]\to\R$ such that
\[
\int_a^x g(t)\,\dd t=f(x)-f(a)
\]
for all $x\in[a,b]$, show that $f$ is continuously differentiable.
\end{problemsubparts}
\end{problem}

% 2022 Spring POSTECH Calculus I Final, Problem 6
\begin{problem}[2022 Spring \POSTECH{} Calculus I]
Let $f$ be a twice continuously differentiable function on an interval $I$ containing $a$ and $b$. Show that
\[
\left|
\frac12\bigl(f(a)+f(b)\bigr)-f\!\left(\frac{a+b}{2}\right)
\right|
\le
\frac18 M(b-a)^2,
\]
where
\[
M=\max_{x\in I}|f''(x)|.
\]

\textit{Hint.} Use the Taylor expansion of the function at the two points $a$ and $b$.
\end{problem}

% 2022 Spring POSTECH Calculus I Final, Problem 7
\begin{problem}[2022 Spring \POSTECH{} Calculus I]
Let
\[
f(x)=\int_{\cos x}^{\sin x}\sin(\cos t)\,\dd t.
\]
Find $f'(x)$.
\end{problem}

% 2022 Spring POSTECH Calculus I Final, Problem 8
\begin{problem}[2022 Spring \POSTECH{} Calculus I]
Let
\[
f(x)=x^6-3x^4+2x^3-x+1.
\]
Find the Taylor polynomial $t_3(x)$ of degree $3$ of $f$ at $1$, and show that
\[
|f(0.3)-t_3(0.3)|\le 12(0.7)^4.
\]
\end{problem}

% 2022 Spring POSTECH Calculus I Final, Problem 9
\begin{problem}[2022 Spring \POSTECH{} Calculus I]
Prove that, for $x>0$,
\[
\frac{x}{1+x^2}
\le
\int_0^x\frac{1}{1+t^2}\,\dd t
\le x.
\]

\textit{Hint.} Use the mean value theorem for integrals.
\end{problem}

% 2022 Spring POSTECH Calculus I Final, Problem 10
\begin{problem}[2022 Spring \POSTECH{} Calculus I]
Show that
\[
f(x)=\frac{1}{1+x^2}
\]
is uniformly continuous on $(-\infty,\infty)$.
\end{problem}

% 2022 Spring POSTECH Calculus I Final, Problem 11
\begin{problem}[2022 Spring \POSTECH{} Calculus I]
Let $f$ be a three-times continuously differentiable function on the interval $[0,2]$.
\begin{problemsubparts}
\item Derive the following equality for any $x\in[0,1]$:
\[
f(x)
=
f(1)+f'(1)(x-1)
+\frac12 f''(1)(x-1)^2
+\frac12\int_1^x (x-t)^2f'''(t)\,\dd t.
\]

\textit{Hint.} Use the fundamental theorem of calculus.
\item Using the intermediate value theorem, show that there is a number $c$ between $1$ and $x$ such that
\[
\int_1^x (x-t)^2f'''(t)\,\dd t
=
\frac13 f'''(c)(x-1)^3.
\]

\textit{Hint.} There exist $m$ and $M$ such that
\[
m\le f'''(x)\le M.
\]
\end{problemsubparts}
\end{problem}


% 2022 Spring KAIST Calculus I Midterm, Problem 1
\begin{problem}[2022 Spring \KAIST{} Calculus I]
Use the precise $\varepsilon$-$\delta$ definitions to prove the following.
\begin{problemsubparts}
\item
\[
\lim_{x\to0}x^2=0.
\]
\item The function
\[
f(x)=\sqrt{|x|}
\]
is continuous at $x=0$.
\end{problemsubparts}
\end{problem}

% 2022 Spring KAIST Calculus I Midterm, Problem 2
\begin{problem}[2022 Spring \KAIST{} Calculus I]
\begin{problemsubparts}
\item Let $(a,b)$ be a point on the circle
\[
x^2+y^2=25.
\]
Show that the tangent line at $(a,b)$ is perpendicular to the line joining $(a,b)$ to the origin.
\item The radius of a sphere is measured as $x$ with an error of at most $1\%$. Using estimates of percentage change, find the maximum percentage errors in the volume and surface area of the sphere.
\end{problemsubparts}
\end{problem}

% 2022 Spring KAIST Calculus I Midterm, Problem 3
\begin{problem}[2022 Spring \KAIST{} Calculus I]
Let
\[
f(x)=x^{1/x},
\qquad x>0.
\]
\begin{problemsubparts}
\item Find
\[
\lim_{x\to\infty}f(x)
\qquad\text{and}\qquad
\lim_{x\to0^+}f(x).
\]
\item Find $f'(x)$ and determine the zero of the derivative, i.e., solve $f'(x)=0$.
\item Sketch the graph of $f$ using the information above.
\end{problemsubparts}
\end{problem}

% 2022 Spring KAIST Calculus I Midterm, Problem 4
\begin{problem}[2022 Spring \KAIST{} Calculus I]
\begin{problemsubparts}
\item Which of the following functions are differentiable at $0$? Justify your answer.
\begin{enumerate}[label=\textup{(\roman*)},leftmargin=2.55em,labelsep=0.62em,itemsep=0.18em,topsep=0.42em,parsep=0pt,partopsep=0pt]
\item $|\sin^3 x|$
\item $\sqrt{x+1}$
\item $\sqrt{x^2+|x|}$
\item $\sqrt{x^2+|x|^3}$
\end{enumerate}
\item Show that
\[
\sinh(x+y)=\sinh x\cosh y+\cosh x\sinh y.
\]
\end{problemsubparts}
\end{problem}

% 2022 Spring KAIST Calculus I Midterm, Problem 5
\begin{problem}[2022 Spring \KAIST{} Calculus I]
\begin{problemsubparts}
\item Let $f\colon[a,b]\to[c,d]$ be an increasing one-to-one function with range $[c,d]$, and let $g\colon[c,d]\to[a,b]$ be its inverse. Find
\[
\int_a^b f(x)\,\dd x+\int_c^d g(y)\,\dd y.
\]
\item Let $f$ be a function defined on $\R$ by
\[
f(x)=
\begin{cases}
\arctan(ax^2+bx+2), & x\ge \dfrac{\pi}{2},\\[0.6em]
\dfrac{\pi}{4}\sin x, & -\dfrac{\pi}{2}<x<\dfrac{\pi}{2},\\[0.8em]
\arctan(dx^2+ex-3), & x\le-\dfrac{\pi}{2}.
\end{cases}
\]
Find values of $a,b,d,e$ which ensure that $f$ is differentiable and has an inverse.
\end{problemsubparts}
\end{problem}

% 2022 Spring KAIST Calculus I Midterm, Problem 6
\begin{problem}[2022 Spring \KAIST{} Calculus I]
\begin{problemsubparts}
\item Find
\[
\lim_{x\to0^+}\frac{\ln x}{1/\sin x}.
\]
\item Find
\[
\lim_{x\to0}
\frac{\displaystyle\int_0^{x^2}\frac{\sin t}{t}\,\dd t}{x^2}.
\]
\end{problemsubparts}
\end{problem}

% 2022 Spring KAIST Calculus I Midterm, Problem 7
\begin{problem}[2022 Spring \KAIST{} Calculus I]
Evaluate the following integrals.
\begin{problemsubparts}
\item
\[
\int\frac{1}{x^2+x+1}\,\dd x.
\]
\item
\[
\int\arccosh x\,\dd x.
\]
\end{problemsubparts}
\end{problem}

% 2022 Spring KAIST Calculus I Midterm, Problem 8
\begin{problem}[2022 Spring \KAIST{} Calculus I]
Evaluate the following integrals.
\begin{problemsubparts}
\item
\[
\int_0^{\pi/2}\cos x\,\sinh(\sin x)\,\dd x.
\]
\item
\[
\int_0^{\pi/2}
\left[
\arctan\!\bigl(\arctanh(\sin x)\bigr)
+\arcsin\!\bigl(\tanh(\tan x)\bigr)
\right]\,\dd x.
\]
\end{problemsubparts}
\end{problem}

% 2022 Spring KAIST Calculus I Midterm, Problem 9
\begin{problem}[2022 Spring \KAIST{} Calculus I]
\begin{problemsubparts}
\item Determine the values of $a\in\R$ for which the improper integral
\[
\int_1^\infty \cos(e^{ax})\,\dd x
\]
converges. Justify your answer.
\item Determine whether
\[
\int_0^\infty\frac{\cos^3 x}{\sqrt{x}}\,\dd x
\]
converges. Justify your answer.
\end{problemsubparts}
\end{problem}

% 2022 Spring KAIST Calculus I Midterm, Problem 10
\begin{problem}[2022 Spring \KAIST{} Calculus I]
\begin{problemsubparts}
\item Find the function $y=y(x)$ satisfying
\[
y^2\frac{\dd y}{\dd x}=3x^2y^3-6x^2,
\qquad
y(0)=1.
\]
\item Suppose the half-life of a radioactive material $X$ is $6$ months. Let $y(t)$ be the mass of $X$ at time $t$, where time is measured in months. Find the decay rate $k$ and the differential equation satisfied by $y$.
\end{problemsubparts}
\end{problem}

% 2022 Spring KAIST Calculus I Final, Problem 1
\begin{problem}[2022 Spring \KAIST{} Calculus I]
Let $f$ satisfy
\[
x\ln x\,f'(x)+f(x)-x=0,
\qquad
f(e^{-1})=-e^{-1},
\qquad
0<x<1.
\]
\begin{problemsubparts}
\item Find $f$ by solving the differential equation.
\item Define
\[
a_{i+1}=|f(a_i)|,
\qquad
a_1=a\in(0,e^{-1}).
\]
Determine whether
\[
\sum_{i=1}^{\infty}a_i
\]
diverges, converges conditionally, or converges absolutely.
\end{problemsubparts}
\end{problem}
% 2022 Spring KAIST Calculus I Final, Problem 3
\begin{problem}[2022 Spring \KAIST{} Calculus I]
\begin{problemsubparts}
\item Suggest three different ways to find the Maclaurin series for $\cos^2 x$. Then use one of them to find the first three nonzero terms of the series.
\item Compute
\[
\lim_{t\to0^+}
\frac{
 e^{-1}\left(1+\frac{t}{2}\right)
 -\left(\frac{1}{1+t}\right)^{1/t}
}{t^2}.
\]
\end{problemsubparts}
\end{problem}
% 2022 Spring KAIST Calculus I Final, Problem 4
\begin{problem}[2022 Spring \KAIST{} Calculus I]
Consider the series
\[
f(x)=1+\sum_{n=1}^{\infty}\frac{(2n)!}{(n!)^2}x^n.
\]
\begin{problemsubparts}
\item Find the radius $R$ of convergence and show that $f$ satisfies
\[
(1-4x)f'(x)=2f(x)
\]
for $x\in(-R,R)$.
\item Find $f(x)$ and compute
\[
1+\sum_{n=1}^{\infty}\frac{(2n)!}{(n!)^2}\frac{1}{8^n}.
\]
\end{problemsubparts}
\end{problem}
% 2022 Spring KAIST Calculus I Final, Problem 5
\begin{problem}[2022 Spring \KAIST{} Calculus I]
Determine whether each of the following sequences converges. Find its limit if it converges. You do not need to justify your answers.
\begin{problemsubparts}
\item $\displaystyle a_n=\left(1+\frac{2}{n}\right)^n$.
\item $\displaystyle b_n=\left(1+\frac{1}{n}\right)^{n^2}$.
\item $\displaystyle c_n=\arctan(\sqrt n-\sin n)$.
\item $\displaystyle d_{n+1}=\frac{4}{3+d_n}$ with $d_1=10$.
\end{problemsubparts}
\end{problem}
% 2022 Spring KAIST Calculus I Final, Problem 7
\begin{problem}[2022 Spring \KAIST{} Calculus I]
For what values of $x$ does each of the following series converge?
\begin{problemsubparts}
\item
\[
\sum_{n=1}^{\infty}\frac{(2x-1)^n}{n}.
\]
\item
\[
\sum_{n=1}^{\infty}\frac{(3x+2)^{n^2}}{n^2}.
\]
\end{problemsubparts}
\end{problem}
% 2022 Spring KAIST Calculus I Final, Problem 9
\begin{problem}[2022 Spring \KAIST{} Calculus I]
\begin{problemsubparts}
\item Show that
\[
\ln(n!)\ge n\ln n-n
\]
for every positive integer $n$.

\textit{Hint.} Use an appropriate Riemann sum of the function $y=\ln x$ on the interval $[1,n]$.
\item Determine whether
\[
\sum_{n=2}^{\infty}(-1)^n\frac{1}{\ln(n!)}
\]
diverges, converges absolutely, or converges conditionally.
\item Determine whether
\[
\sum_{n=2}^{\infty}(-1)^n
\frac{1}{\ln^2(n)\sqrt[n]{n!}}
\]
diverges, converges absolutely, or converges conditionally.
\end{problemsubparts}
\end{problem}

% 2023 Spring POSTECH Calculus I Midterm, Problem 1
\begin{problem}[2023 Spring \POSTECH{} Calculus I]
Prove the following statements.
\begin{problemsubparts}
\item The limit of a convergent sequence is unique.
\item Every convergent sequence is a Cauchy sequence.
\end{problemsubparts}
\end{problem}

% 2023 Spring POSTECH Calculus I Midterm, Problem 2
\begin{problem}[2023 Spring \POSTECH{} Calculus I]
Consider the sequence $(a_n)$ defined by
\[
a_1=2,
\qquad
a_{n+1}=\frac12\left(a_n+\frac{3}{a_n}\right)
\quad\text{for all integers }n>1.
\]
Prove the following statements.
\begin{problemsubparts}
\item $a_n>\sqrt3$ for all integers $n\ge1$.
\item The sequence $(a_n)$ is decreasing; that is, $a_n>a_{n+1}$.
\item $\displaystyle\lim_{n\to\infty}a_n=\sqrt3$.
\end{problemsubparts}
\end{problem}

% 2023 Spring POSTECH Calculus I Midterm, Problem 3
\begin{problem}[2023 Spring \POSTECH{} Calculus I]
Let $f\colon[0,1]\to[0,1]$ be continuous. Show that there exists $c\in[0,1]$ such that
\[
f(c)=c.
\]
\end{problem}

% 2023 Spring POSTECH Calculus I Midterm, Problem 4
\begin{problem}[2023 Spring \POSTECH{} Calculus I]
Let $a<b$, and let $f\colon[a,b]\to\R$ be continuous. Show that the range of $f$ is a closed interval.
\end{problem}

% 2023 Spring POSTECH Calculus I Midterm, Problem 5
\begin{problem}[2023 Spring \POSTECH{} Calculus I]
Let $a$ be any positive real number. Show that
\[
f(x)=x^{1/3}
\]
is uniformly continuous on $[a,\infty)$.
\end{problem}

% 2023 Spring POSTECH Calculus I Midterm, Problem 6
\begin{problem}[2023 Spring \POSTECH{} Calculus I]
Consider the function
\[
f(x)=(x^3+1)^{1/2}
\]
on the interval $[-1,3]$.
\begin{problemsubparts}
\item Sketch the graph of $f$.
\item Show that $f$ is strictly monotone on $[-1,3]$ and find the inverse function $g=f^{-1}$. Give the domain and range of $g$.
\item Is the inverse $g$ differentiable on its domain?
\end{problemsubparts}
\end{problem}

% 2023 Spring POSTECH Calculus I Midterm, Problem 7
\begin{problem}[2023 Spring \POSTECH{} Calculus I]
Let
\[
\sum_{n=1}^{\infty}a_n
\]
be a series of positive terms. If this series converges, show that
\[
\sum_{n=1}^{\infty}\ln(1+a_n)
\]
converges.

\textit{Hint.} Use the limit comparison test.
\end{problem}

% 2023 Spring POSTECH Calculus I Midterm, Problem 8
\begin{problem}[2023 Spring \POSTECH{} Calculus I]
Consider the second-order differential equation
\[
y''(t)+4y(t)=0
\]
for all $t\in(-\infty,\infty)$, with the conditions
\[
y(0)=1,
\qquad
y'(0)=1.
\]
\begin{problemsubparts}
\item Find a solution $y(t)$.
\item Show that the solution is unique.
\end{problemsubparts}
\end{problem}

% 2023 Spring POSTECH Calculus I Midterm, Problem 9
\begin{problem}[2023 Spring \POSTECH{} Calculus I]
Let
\[
f(x)=
\begin{cases}
x\sin\!\left(\dfrac{\pi}{x}\right), & x\ne0,\\[0.5em]
0, & x=0,
\end{cases}
\qquad
\text{and}
\qquad
g(x)=
\begin{cases}
x^2\sin\!\left(\dfrac{\pi}{x}\right), & x\ne0,\\[0.5em]
0, & x=0.
\end{cases}
\]
\begin{problemsubparts}
\item Determine whether $f$ is differentiable at $0$.
\item Determine whether $g$ is differentiable at $0$.
\end{problemsubparts}
\end{problem}

% 2023 Spring POSTECH Calculus I Midterm, Problem 10
\begin{problem}[2023 Spring \POSTECH{} Calculus I]
Consider the power series
\[
f(x)=\sum_{n=1}^{\infty}\frac{x^n}{n}.
\]
\begin{problemsubparts}
\item Determine all $x$ for which $f(x)$ converges.
\item Evaluate $f''(1/2)$.
\end{problemsubparts}
\end{problem}


% 2023 Spring KAIST Calculus I Midterm, Problem 1
\begin{problem}[2023 Spring \KAIST{} Calculus I]
Use the precise $\varepsilon$-$\delta$ definitions for the following.
\begin{problemsubparts}
\item Suppose
\[
\lim_{x\to c}f(x)=L
\qquad\text{and}\qquad
\lim_{x\to c}g(x)=M,
\]
where $L,M,c\in\R$. Prove that
\[
\lim_{x\to c}\bigl(f(x)+g(x)\bigr)=L+M.
\]
\item Determine whether
\[
f(x)=
\begin{cases}
x^2, & x\ne0,\\
2, & x=0
\end{cases}
\]
is continuous at $x=0$.
\end{problemsubparts}
\end{problem}

% 2023 Spring KAIST Calculus I Midterm, Problem 2
\begin{problem}[2023 Spring \KAIST{} Calculus I]
\begin{problemsubparts}
\item Find all horizontal asymptotes of the graph of
\[
f(x)=\frac{x^3-|x|^3+x^2}{|x|^3+2},
\]
if any.
\item Prove that
\[
e^x-3x^2=0
\]
has at least three solutions. Recall that $e\approx2.71828$.
\end{problemsubparts}
\end{problem}

% 2023 Spring KAIST Calculus I Midterm, Problem 3
\begin{problem}[2023 Spring \KAIST{} Calculus I]
Consider
\[
f(x)=
\begin{cases}
x^2\sin(1/x), & x\ne0,\\
0, & x=0.
\end{cases}
\]
\begin{problemsubparts}
\item Show that $f$ is continuous at $x=0$.
\item Show that $f'(x)$ exists for $x\ne0$.
\item Show that $f$ is differentiable at $x=0$.
\item Show that $f'$ is not continuous at $x=0$.
\end{problemsubparts}
\end{problem}

% 2023 Spring KAIST Calculus I Midterm, Problem 4
\begin{problem}[2023 Spring \KAIST{} Calculus I]
Evaluate
\[
\int_0^\pi \cos(\arccot x)\,\dd x.
\]
\end{problem}

% 2023 Spring KAIST Calculus I Midterm, Problem 5
\begin{problem}[2023 Spring \KAIST{} Calculus I]
Consider the curve
\[
y^4=y^2-x^2
\]
in the $xy$-plane.
\begin{problemsubparts}
\item Find the equation of the tangent line at
\[
\left(\frac{\sqrt3}{4},\frac{\sqrt3}{2}\right).
\]
\item Find the equation of the normal line at
\[
\left(\frac{\sqrt3}{4},\frac12\right).
\]
\end{problemsubparts}
\end{problem}

% 2023 Spring KAIST Calculus I Midterm, Problem 6
\begin{problem}[2023 Spring \KAIST{} Calculus I]
Let
\[
f(x)=\sqrt{1+x^3}=(1+x^3)^{1/2}.
\]
\begin{problemsubparts}
\item Find a polynomial that approximates $f(x)$ at $x=0$, using
\[
(1+x)^k\approx1+kx
\]
for small $x$.
\item Use the product rule to find $f(0)$, $f'(0)$, $f''(0)$, and $f'''(0)$, and then find
\[
p(x)=f(0)+f'(0)x+\frac{f''(0)}{2!}x^2+\frac{f'''(0)}{3!}x^3.
\]
\item Find the linearization at $x=0$ of
\[
g(x)=\int_3^{3+\sin x}f(t)\,\dd t.
\]
\end{problemsubparts}
\end{problem}

% 2023 Spring KAIST Calculus I Midterm, Problem 7
\begin{problem}[2023 Spring \KAIST{} Calculus I]
Evaluate the integral and express the answer in terms of the natural logarithm:
\[
\int_1^{\sqrt5}2x\,\arccosh(x^2)\,\dd x.
\]
\end{problem}

% 2023 Spring KAIST Calculus I Midterm, Problem 8
\begin{problem}[2023 Spring \KAIST{} Calculus I]
Evaluate
\[
\int_0^\infty\frac{1}{\sqrt{x(x+4)}}\,\dd x
\]
by splitting it into the following two integrals.
\begin{problemsubparts}
\item
\[
\int_0^4\frac{1}{\sqrt{x(x+4)}}\,\dd x.
\]
\item
\[
\int_4^\infty\frac{1}{\sqrt{x(x+4)}}\,\dd x.
\]
\end{problemsubparts}
\end{problem}

% 2023 Spring KAIST Calculus I Midterm, Problem 9
\begin{problem}[2023 Spring \KAIST{} Calculus I]
Let $y$ be a solution of
\[
y'=-\frac{2xy}{1+x^2},
\qquad
y(0)=1.
\]
\begin{problemsubparts}
\item Find $y$.
\item Find
\[
\int_{-\infty}^{\infty}y(x)\,\dd x.
\]
\end{problemsubparts}
\end{problem}

% 2023 Spring KAIST Calculus I Midterm, Problem 10
\begin{problem}[2023 Spring \KAIST{} Calculus I]
Let $a>0$ and $b>0$ be constants satisfying $ab^2>1$. Let
\[
I=[2,\infty),
\]
and let $f$ be differentiable on $I$ with $f(2)=b^2$. Suppose that, with $y=f(x)$,
\[
2axy\frac{\dd y}{\dd x}-2y=x\frac{\dd y}{\dd x}.
\]
Show that $f$ is one-to-one on $I$, and find $f^{-1}$.
\end{problem}


% 2023 Spring KAIST Calculus I Final, Problem 1
\begin{problem}[2023 Spring \KAIST{} Calculus I]
Let
\[
f(x)=\arctan x,
\qquad -\infty<x<\infty.
\]
\begin{problemsubparts}
\item Find the set of points at which the Maclaurin series generated by $f$ converges to $f$.
\item Use the Maclaurin series of $\arctan x$ to show that
\[
\left|
\frac{\pi}{\sqrt3}
-\left(2-\frac29+\frac{2}{45}-\frac{2}{189}\right)
\right|
\le\frac{2}{3^6}.
\]
\end{problemsubparts}
\end{problem}

% 2023 Spring KAIST Calculus I Final, Problem 2
\begin{problem}[2023 Spring \KAIST{} Calculus I]
Determine whether each of the following series converges.
\begin{problemsubparts}
\item
\[
\sum_{n=2}^{\infty}\frac{(-1)^n n}{(n+1)\ln n}.
\]
\item
\[
\sum_{n=1}^{\infty}(-1)^n a_n^{a_n},
\]
where $a_n>0$ for all $n$ and
\[
\lim_{n\to\infty}a_n=\frac1e\approx\frac1{2.71828}.
\]
\end{problemsubparts}
\end{problem}

% 2023 Spring KAIST Calculus I Final, Problem 3
\begin{problem}[2023 Spring \KAIST{} Calculus I]
A tank contains $200\,\mathrm{L}$ of water with $10\,\mathrm{kg}$ of salt dissolved in it. Salt water containing $0.2\,\mathrm{kg}$ of salt per liter is added at a rate of $20\,\mathrm{L/min}$, while the well-mixed solution is pumped out at a rate of $10\,\mathrm{L/min}$.
\begin{problemsubparts}
\item Let $y(t)$ be the kilograms of salt in the tank $t$ minutes after the pumping process begins. Write a first-order linear differential equation satisfied by $y(t)$.
\item How much salt is in the tank $4$ minutes after the pumping process begins?
\end{problemsubparts}
\end{problem}

% 2023 Spring KAIST Calculus I Final, Problem 9
\begin{problem}[2023 Spring \KAIST{} Calculus I]
\begin{problemsubparts}
\item Find the interval of convergence of
\[
\sum_{n=1}^{\infty}\frac{(-1)^n x^n}{\sqrt{n+1}}.
\]
\item Find the interval of convergence of
\[
\sum_{n=1}^{\infty}\left(\frac{x^3+1}{2}\right)^n
\]
and find the sum of the series as a simple function of $x$.
\item Find the Maclaurin series of the function obtained in part \textup{(b)}, together with its interval of convergence.
\end{problemsubparts}
\end{problem}

% 2023 Spring KAIST Calculus I Final, Problem 10
\begin{problem}[2023 Spring \KAIST{} Calculus I]
\begin{problemsubparts}
\item Show that
\[
\ln(n+1)
\le
1+\frac12+\cdots+\frac1n
\le
1+\ln n.
\]
\item Show that
\[
\frac1{n+1}<\int_n^{n+1}\frac1x\,\dd x.
\]
\item Let
\[
a_n=1+\frac12+\cdots+\frac1n-\ln n.
\]
Show that the sequence $(a_n)$ converges as $n\to\infty$.
\end{problemsubparts}
\end{problem}


% 2024 Spring KAIST Calculus I Midterm, Problem 1
\begin{problem}[2024 Spring \KAIST{} Calculus I]
\begin{problemsubparts}
\item Let
\[
f(x)=x^2+3.
\]
Using the $\varepsilon$-$\delta$ definition of continuity, show that $f$ is continuous at $x=1$.
\item Use the mean value theorem to show that, for any distinct $\theta_1,\theta_2\in[0,\pi/4]$,
\[
\left|
\frac{\ln(\cos\theta_1)-\ln(\cos\theta_2)}{\theta_1-\theta_2}
\right|
\le1.
\]
\end{problemsubparts}
\end{problem}

% 2024 Spring KAIST Calculus I Midterm, Problem 2
\begin{problem}[2024 Spring \KAIST{} Calculus I]
Suppose $x,y,z$ satisfy
\[
x^3+xy+y^3=1,
\qquad
\cos(\pi y)+\sin(\pi y)=z^3.
\]
\begin{problemsubparts}
\item Find $\dd y/\dd x$ at $x=0$.
\item Find $\dd z/\dd x$ at $x=0$.
\end{problemsubparts}
\end{problem}

% 2024 Spring KAIST Calculus I Midterm, Problem 3
\begin{problem}[2024 Spring \KAIST{} Calculus I]
Let
\[
f(x)=-x^3-4x^2-x+3.
\]
\begin{problemsubparts}
\item Find the largest interval $[a,b]$ containing $-1$ on which $f$ is one-to-one.
\item Let $g$ be the inverse of the restriction of $f$ to $[a,b]$. Find $g'(-3)$ if it exists, or prove that it does not exist.
\end{problemsubparts}
\end{problem}

% 2024 Spring KAIST Calculus I Midterm, Problem 4
\begin{problem}[2024 Spring \KAIST{} Calculus I]
\begin{problemsubparts}
\item Determine whether
\[
f(x)=
\begin{cases}
x^2\sin(\ln|x|), & x\ne0,\\
0, & x=0
\end{cases}
\]
is differentiable at $x=0$.
\item Find the derivative of
\[
y=(1+x^2)^{\sin x}.
\]
\end{problemsubparts}
\end{problem}

% 2024 Spring KAIST Calculus I Midterm, Problem 6
\begin{problem}[2024 Spring \KAIST{} Calculus I]
A deposit is left in a bank account for $10$ years under compound interest, and the total amount doubles. Construct a differential equation for the amount and find the annual interest rate.
\end{problem}

% 2024 Spring KAIST Calculus I Midterm, Problem 7
\begin{problem}[2024 Spring \KAIST{} Calculus I]
Consider the differential equation
\[
y''-2y'+y=x^2+x.
\]
\begin{problemsubparts}
\item Find a particular solution $y_p$. \textit{Hint.} $y_p=x^2+ax+b$ for some constants $a,b$.
\item Find the unique solution satisfying
\[
y'(0)=0,
\qquad
y(0)=0.
\]
\end{problemsubparts}
\end{problem}


% 2024 Spring KAIST Calculus I Final, Problem 2
\begin{problem}[2024 Spring \KAIST{} Calculus I]
Determine whether each of the following limits exists, and if it does, find the limit.
\begin{problemsubparts}
\item
\[
\lim_{x\to0}\frac{x-\sin x}{\sinh x-x}.
\]
\item
\[
\lim_{x\to0}
\frac{\displaystyle\int_0^{x^2}\arctan t\,\dd t}{x^4}.
\]
\item
\[
\lim_{x\to\infty}\bigl[\ln(e^x+1)-x\bigr].
\]
\item
\[
\lim_{x\to\infty}x\left(\frac\pi2-\arctan x\right).
\]
\end{problemsubparts}
\end{problem}

% 2024 Spring KAIST Calculus I Final, Problem 3
\begin{problem}[2024 Spring \KAIST{} Calculus I]
Let $f$ be continuously differentiable and one-to-one on $[a,b]$, and assume $f'(x)\ne0$ on $[a,b]$. Compute
\[
\int_{f(a)}^{f(b)}f^{-1}(y)\,\dd y
+
\int_a^b f(x)\,\dd x.
\]
\end{problem}

% 2024 Spring KAIST Calculus I Final, Problem 4
\begin{problem}[2024 Spring \KAIST{} Calculus I]
Let
\[
\Gamma(n)=\int_0^\infty x^{n-1}e^{-x}\,\dd x,
\qquad n\ge1.
\]
\begin{problemsubparts}
\item Show by the comparison test that the improper integral $\Gamma(n)$ converges for every real $n\ge1$.
\item Find $\Gamma(n)$ when $n$ is a positive integer.
\end{problemsubparts}
\end{problem}

% 2024 Spring KAIST Calculus I Final, Problem 5
\begin{problem}[2024 Spring \KAIST{} Calculus I]
Evaluate
\[
\int_{-\infty}^{0}
\frac{7x^2-2}{(x-1)^2(x^2+2x+2)}\,\dd x.
\]
\end{problem}

% 2024 Spring KAIST Calculus I Final, Problem 6
\begin{problem}[2024 Spring \KAIST{} Calculus I]
Test the convergence of the following series.
\begin{problemsubparts}
\item
\[
\sum_{n=1}^{\infty}\frac{(n!)^2 4^n}{(2n)!}.
\]
\item
\[
\sum_{n=1}^{\infty}
\frac{(-1)^n4^n+n^{617}3^n}{n4^n}.
\]
\end{problemsubparts}
\end{problem}

% 2024 Spring KAIST Calculus I Final, Problem 7
\begin{problem}[2024 Spring \KAIST{} Calculus I]
Prove or disprove the following statements.
\begin{problemsubparts}
\item If
\[
\sum_{n=1}^{\infty}a_n
\qquad\text{and}\qquad
\sum_{n=1}^{\infty}b_n
\]
converge, then
\[
\sum_{n=1}^{\infty}a_nb_n
\]
converges.
\item If the two series above converge absolutely, then
\[
\sum_{n=1}^{\infty}a_nb_n
\]
converges absolutely.
\end{problemsubparts}
\end{problem}

% 2024 Spring KAIST Calculus I Final, Problem 8
\begin{problem}[2024 Spring \KAIST{} Calculus I]
Determine the values of $x$ for which the following series converge.
\begin{problemsubparts}
\item
\[
\sum_{n=2}^{\infty}\frac{(-x)^n}{n(\ln n)^2}.
\]
\item
\[
\sum_{n=1}^{\infty}n(\sin x)^{n-1}.
\]
\end{problemsubparts}
\end{problem}

% 2024 Spring KAIST Calculus I Final, Problem 9
\begin{problem}[2024 Spring \KAIST{} Calculus I]
\begin{problemsubparts}
\item Find the Taylor polynomial $p_n(x)$ for $\sin x$ centered at $x_0=\pi/2$, and state Taylor's theorem using this polynomial.
\item Use Taylor's theorem to show that
\[
p_n(x)\to\sin x
\]
for every $x\in\R$ as $n\to\infty$.
\end{problemsubparts}
\end{problem}

% 2024 Spring KAIST Calculus I Final, Problem 10
\begin{problem}[2024 Spring \KAIST{} Calculus I]
\begin{problemsubparts}
\item Use Taylor's theorem to find the second-order term of
\[
(2+x+3x^2)^{10}.
\]
\item Let
\[
f(x)=\sqrt{\arctan x+1}.
\]
Approximate $f(2)$ using the linear approximation of $f$ at $x=1$.
\end{problemsubparts}
\end{problem}

